% ORBIT-MS WGV / RSI supplementary results.
% Added from orbit-ms PR95 (head 6a5fa7c, 2026-11-09) and its linked PR93
% result artifacts. Tables S14 onward follow existing Tables S1--S13.
% Important: cohorts and task definitions differ across tables; do not compare
% their percentages as if they came from one common evaluation protocol.

\clearpage
\section{WGV generation, recursive improvement and candidate verification}
\label{sec:wgv_rsi_results}

This section consolidates the WGV and recursive self-improvement (RSI) experiments
reported in ORBIT-MS PR95 and the associated PR93 result artifacts. The experiments
cover distinct evaluation settings: held-out Calibration with the reference formula
provided; official-test development snapshots with predicted formula hypotheses;
TRAIN-derived diagnostic subsets; and database-assisted ranking from a fixed
candidate pool. Exact connectivity is assessed using the first 14 characters of the
InChIKey and does not establish stereochemical identity. \emph{Reach} denotes the
presence of the reference connectivity in the generated candidate set, whereas
Tanimoto scores are post-hoc structural-similarity measures.

The generation ledger in Table~\ref{tab:rsi_gen_ledger} records successive model and
inference updates. It is a chronological development trajectory, not a record of
seven independently completed autonomous-agent episodes. The later Gen~6 and Gen~7
entries are interim partial-evaluation snapshots. The official TEST subset was
consulted during iterative development; these results should therefore be described
as an official-test development ledger rather than as a final untouched blind test.
The complete source record is available in
\url{https://github.com/wangyu-sd/orbit-ms/pull/95}.

\begin{table}[!htbp]
\centering
\caption{\textbf{Official-test development trajectory across WGV generations.}
Successive changes include HMDB-W training, unified MMA expansion, atom-normalized
search, macro-core initialization, short-horizon G completion and second-round
candidate-guided W/V refinement. Gen~6 and Gen~7 are interim partial snapshots
(marked by $\dagger$), not final full-cohort estimates.}
\label{tab:rsi_gen_ledger}
\small
\setlength{\tabcolsep}{3.5pt}
\renewcommand{\arraystretch}{1.10}
\noindent\textit{A. Connectivity recovery}\par\vspace{3pt}
\begin{tabular}{@{}lp{0.49\linewidth}rr@{}}
\toprule
Generation & Main intervention & Exact@1 & Exact@10 \\
\midrule
Gen 1 & TRAIN-only G; MMA library v1 & 0.00\% & 1.00\% \\
Gen 2 & HMDB-W training; closed-motif sets & 0.00\% & 0.00\% \\
Gen 3 & Atom-normalized beam search & 0.50\% & 0.50\% \\
Gen 4 & Expanded 12.98M-entry MMA library & 2.00\% & 2.00\% \\
Gen 5 & Dual starts; first-order post-W/V refinement & 2.00\% & 2.50\% \\
Gen 6$^\dagger$ & Non-TEST macro-core pruning; short-horizon G completion & 4.00\% & 10.00\% \\
Gen 7$^\dagger$ & Isomer consensus; topology prior; second-round W/V & \textbf{10.00\%} & \textbf{10.00\%} \\
\bottomrule
\end{tabular}
\par\vspace{8pt}
\noindent\textit{B. Structural similarity}\par\vspace{3pt}
\begin{tabular}{@{}lrrr@{}}
\toprule
Generation & Tanimoto@1 & Tanimoto@10 & Pool maximum \\
\midrule
Gen 1 & 0.1612 & 0.2175 & 0.2480 \\
Gen 2 & 0.1699 & 0.2142 & 0.2410 \\
Gen 3 & 0.1799 & 0.2174 & 0.2495 \\
Gen 4 & 0.1904 & 0.2291 & 0.2700 \\
Gen 5 & 0.1972 & 0.2469 & 0.2840 \\
Gen 6$^\dagger$ & 0.2803 & 0.3603 & 0.4114 \\
Gen 7$^\dagger$ & \textbf{0.2904} & \textbf{0.3603} & \textbf{0.4114} \\
\bottomrule
\end{tabular}
\end{table}

\paragraph{Calibration progression.}
The official-test ledger is complemented by the held-out Calibration progression
in Table~\ref{tab:calibration_progression}. These runs use the reference molecular
formula and therefore measure structure generation under a known-formula condition.
They are kept separate from the target-free official-test results above.

\begin{table}[!htbp]
\centering
\caption{\textbf{Held-out Calibration progression from Atom-only generation to Canonical Hybrid WGV.}
All rows use the same Calibration evaluation protocol with the reference formula
provided. Gen labels here denote successive development configurations and are
not the official-test ledger labels in Table~\ref{tab:rsi_gen_ledger}.}
\label{tab:calibration_progression}
\small
\setlength{\tabcolsep}{3.5pt}
\renewcommand{\arraystretch}{1.10}
\noindent\textit{A. Connectivity recovery}\par\vspace{3pt}
\begin{tabular}{@{}lp{0.42\linewidth}rrr@{}}
\toprule
Stage & Configuration & Exact@1 & Exact@10 & Reach \\
\midrule
Gen 0 & Atom-only; no MMA & 1.50\% & 4.50\% & 6.50\% \\
Gen 1 & MMA features & 6.00\% & 14.50\% & 19.00\% \\
Gen 2 & MMA features; W/V terminal ranking & 8.50\% & 13.50\% & 16.00\% \\
Gen 3 & Eight MMA starts; W/V ranking & 10.00\% & 18.00\% & 19.00\% \\
Gen 4 & Canonical Hybrid; atom-normalized and motif-reserved beam & \textbf{12.50\%} & \textbf{21.50\%} & \textbf{22.50\%} \\
\bottomrule
\end{tabular}
\par\vspace{8pt}
\noindent\textit{B. Structural similarity}\par\vspace{3pt}
\begin{tabular}{@{}lrr@{}}
\toprule
Stage & Tanimoto@1 & Tanimoto@10 \\
\midrule
Gen 0 & 0.2009 & 0.2640 \\
Gen 1 & 0.3221 & 0.4215 \\
Gen 2 & 0.3351 & 0.4152 \\
Gen 3 & 0.3752 & 0.4685 \\
Gen 4 & \textbf{0.3811} & \textbf{0.4938} \\
\bottomrule
\end{tabular}
\end{table}

\subsection{MMA knowledge and generation ablations}

\begin{table}[!htbp]
\centering
\caption{\textbf{Strict same-checkpoint ablation of MMA knowledge.}
The same full-TRAIN checkpoint and decoding configuration are evaluated with
MMA enabled or disabled. Teacher and action-level metrics assess the ladder
decisions; Exact@$k$, Reach and Tanimoto assess complete candidate generation.}
\label{tab:mma_strict_ablation}
\small
\setlength{\tabcolsep}{3.5pt}
\renewcommand{\arraystretch}{1.10}
\noindent\textit{A. Teacher and action accuracy}\par\vspace{3pt}
\begin{tabular}{@{}lrrrr@{}}
\toprule
Configuration & Teacher loss $\downarrow$ & Teacher@1 & Start@1 & Select@1 \\
\midrule
Without MMA & 0.996 & 69.1\% & 76.4\% & 54.6\% \\
With MMA & \textbf{0.757} & \textbf{76.0\%} & \textbf{86.3\%} & \textbf{67.5\%} \\
\bottomrule
\end{tabular}
\par\vspace{8pt}
\noindent\textit{B. Full-molecule generation}\par\vspace{3pt}
\begin{tabular}{@{}lrrrrr@{}}
\toprule
Configuration & Exact@1 & Exact@10 & Reach & Tanimoto@1 & Tanimoto@10 \\
\midrule
Without MMA & 1.50\% & 4.50\% & 6.50\% & 0.2009 & 0.2640 \\
With MMA & \textbf{6.00\%} & \textbf{14.50\%} & \textbf{19.00\%} & \textbf{0.3221} & \textbf{0.4215} \\
\bottomrule
\end{tabular}
\end{table}

\begin{table}[!htbp]
\centering
\caption{\textbf{Contribution of MMA features, MMA-seeded starts and W/V ranking.}
Rows are successive component configurations on the same held-out Calibration
protocol. The table separates candidate generation changes from terminal
verification and ranking changes.}
\label{tab:mma_wv_component_ablation}
\small
\setlength{\tabcolsep}{3.5pt}
\renewcommand{\arraystretch}{1.10}
\noindent\textit{A. Connectivity recovery and candidate-pool reach}\par\vspace{3pt}
\begin{tabular}{@{}lp{0.47\linewidth}rrr@{}}
\toprule
ID & Configuration & Exact@1 & Exact@10 & Reach \\
\midrule
C1 & No MMA; empty-graph start & 1.50\% & 4.50\% & 6.50\% \\
C2 & MMA features; empty start; G ranking & 7.00\% & 11.50\% & 16.00\% \\
C3 & MMA features; empty start; G + W/V & 8.50\% & 13.50\% & 16.00\% \\
C4 & Eight MMA starts; G ranking only & 5.50\% & 16.00\% & 19.00\% \\
C5 & Eight MMA starts; G + W/V & 10.00\% & 18.00\% & 19.00\% \\
C6 & Joint empty-start Beam-16 and MMA-start pool; W/V & 11.00\% & 19.50\% & 22.00\% \\
C7 & Canonical Hybrid; atom-normalized and motif-reserved beam; W/V & \textbf{12.50\%} & \textbf{21.50\%} & \textbf{22.50\%} \\
\bottomrule
\end{tabular}
\par\vspace{8pt}
\noindent\textit{B. Structural similarity}\par\vspace{3pt}
\begin{tabular}{@{}lrr@{}}
\toprule
ID & Tanimoto@1 & Tanimoto@10 \\
\midrule
C1 & 0.2009 & 0.2640 \\
C2 & 0.3170 & 0.4136 \\
C3 & 0.3351 & 0.4136 \\
C4 & 0.3103 & 0.4685 \\
C5 & 0.3752 & 0.4685 \\
C6 & \textbf{0.3872} & 0.4812 \\
C7 & 0.3811 & \textbf{0.4938} \\
\bottomrule
\end{tabular}
\end{table}

\begin{table}[!htbp]
\centering
\caption{\textbf{Terminal W/V reranking on a fixed G-generated candidate pool.}
Because the candidate pool is held fixed, differences isolate the contribution
of the ranking terms rather than changes in proposal generation.}
\label{tab:wv_fixed_pool_reranking}
\small
\setlength{\tabcolsep}{3pt}

\begin{tabular}{@{}p{0.37\linewidth}rrrrr@{}}
\toprule
Ranking strategy & Exact@1 & Exact@10 & Reach & Tan@1 & Tan@10 \\
\midrule
G policy score only & 7.00\% & 11.50\% & 16.00\% & 0.3170 & 0.4136 \\
G + explained intensity & 7.00\% & 11.50\% & 16.00\% & 0.3179 & 0.4136 \\
G + W/V (higher W/V weights) & 7.50\% & 13.50\% & 16.00\% & 0.3327 & 0.4136 \\
G + W/V (selected weights) & \textbf{8.50\%} & \textbf{13.50\%} & 16.00\% & \textbf{0.3351} & 0.4136 \\
\bottomrule
\end{tabular}
\end{table}

\subsection{Canonical Hybrid action space and trajectory efficiency}

\begin{table}[!htbp]
\centering
\caption{\textbf{Canonical Hybrid policy ablation.}
The seen-subset columns report reconstruction of molecules encountered during
the corresponding training regime; the fit-subset Reach column reports
generalization to the designated fit evaluation subset. These endpoints have
different evaluation roles and should not be combined into one accuracy estimate.}
\label{tab:canonical_hybrid_ablation}
\small
\setlength{\tabcolsep}{3.5pt}
\renewcommand{\arraystretch}{1.10}
\noindent\textit{A. Action efficiency and teacher accuracy}\par\vspace{3pt}
\begin{tabular}{@{}lp{0.42\linewidth}rrr@{}}
\toprule
ID & Policy & Decisions & Canonical@1 & Valid@1 \\
\midrule
P1 & Random-order staged baseline, Beam 4 & 32.4 & -- & 90.6\% \\
P2 & Canonical + Motif Hybrid, Beam 4 & 24.5 & 84.7\% & 90.6\% \\
P3 & Direct single-stage generation, Beam 16 & \textbf{9.1} & 74.6\% & 82.9\% \\
P4 & Canonical Hybrid + $R_t$ gate, Beam 16 & 24.6 & \textbf{86.0\%} & \textbf{90.7\%} \\
\bottomrule
\end{tabular}
\par\vspace{8pt}
\noindent\textit{B. Complete-molecule generation}\par\vspace{3pt}
\begin{tabular}{@{}lrrrr@{}}
\toprule
ID & Seen Exact@1 & Seen Reach & Seen Tan@1 & Fit Reach \\
\midrule
P1 & 25.0\% & 37.5\% & 0.488 & 13.3\% \\
P2 & \textbf{32.8\%} & 42.2\% & 0.518 & 14.3\% \\
P3 & 17.2\% & 31.3\% & 0.367 & \textbf{31.3\%} \\
P4 & 28.1\% & \textbf{53.1\%} & \textbf{0.534} & 27.1\% \\
\bottomrule
\end{tabular}
\end{table}

A separate GT-conditioned feasibility probe found that Canonical Hybrid reduced
the mean total decision count from 32.4 to 24.8 and the molecular-growth decision
count from 20.1 to 12.9. Motif actions were used in 87\% of reconstructed molecules,
with a mean of 4.5 heavy atoms attached per Motif action. All probe molecules were
reconstructed under the GT-conditioned canonical action sequence. This is a
trajectory-feasibility result, not an unknown-spectrum structure-identification rate.

\subsection{Database-assisted structure identification}

\begin{table}[!htbp]
\centering
\caption{\textbf{Database-assisted candidate identification with MMA and physical cleavage evidence.}
The blind columns use predicted-formula candidate pools; the oracle-formula columns
provide the reference formula and are a separate, easier condition. This experiment
ranks structures in a fixed database-derived candidate pool and is not a pure
de novo generation result.}
\label{tab:bio_prior_mma_cleavage}
\small
\setlength{\tabcolsep}{3.5pt}
\renewcommand{\arraystretch}{1.10}
\noindent\textit{A. Predicted-formula candidate identification}\par\vspace{3pt}
\begin{tabular}{@{}lp{0.47\linewidth}rrr@{}}
\toprule
ID & Candidate pool and evidence & Exact@1 & Exact@10 & Tan@1 \\
\midrule
D1 & Uniform isomer pool & 17.36\% & 47.02\% & 0.3419 \\
D2 & Uniform pool + MMA alignment & 20.11\% & 50.14\% & 0.3708 \\
D3 & Uniform pool + physical 1/2-cut cleavage & 23.73\% & 52.03\% & 0.3972 \\
D4 & Uniform pool + MMA + physical cleavage & 24.47\% & 53.29\% & 0.4070 \\
D5 & Multi-database Bio-Prior + MMA + physical cleavage & \textbf{31.57\%} & \textbf{56.70\%} & \textbf{0.4583} \\
\bottomrule
\end{tabular}
\par\vspace{8pt}
\noindent\textit{B. Reference-formula candidate identification}\par\vspace{3pt}
\begin{tabular}{@{}lrrr@{}}
\toprule
ID & Exact@1 & Exact@10 & Tan@1 \\
\midrule
D1 & 25.34\% & 54.19\% & 0.4284 \\
D2 & 29.39\% & 57.03\% & 0.4656 \\
D3 & 32.56\% & 58.74\% & 0.4875 \\
D4 & 33.50\% & 59.93\% & 0.4982 \\
D5 & \textbf{42.74\%} & \textbf{63.08\%} & \textbf{0.5625} \\
\bottomrule
\end{tabular}
\end{table}

The source report evaluates this fixed-pool experiment on the complete official
MassSpecGym TEST release (17,556 spectra representing 2,998 unique molecules).
These results measure database-assisted candidate identification, not the
performance of the pure de novo G generator.

\subsection{Representative Gen 7 structure recoveries}

\begin{table}[!htbp]
\centering
\caption{\textbf{Representative exact-connectivity recoveries in the current Gen 7 run.}
The listed cases illustrate large and chemically diverse structures promoted to
rank one by the combined macro-core completion and candidate-ranking pipeline.
The Gen 7 evaluation snapshot remains interim.}
\label{tab:gen7_landmark_cases}
\small
\setlength{\tabcolsep}{2.5pt}

\begin{tabular}{@{}p{0.19\linewidth}p{0.30\linewidth}lrr@{}}
\toprule
MassSpecGym ID & Molecule & Formula & Heavy atoms & Final rank \\
\midrule
0000201 & Beauvericin & C$_{45}$H$_{57}$N$_3$O$_9$ & 57 & 2 $\rightarrow$ 1 \\
0039311 & Roxithromycin & C$_{41}$H$_{76}$N$_2$O$_{15}$ & 58 & 4 $\rightarrow$ 1 \\
0012997 & Clarithromycin & C$_{38}$H$_{69}$NO$_{13}$ & 52 & 5 $\rightarrow$ 1 \\
0012696 & Erythromycin & C$_{37}$H$_{67}$NO$_{13}$ & 51 & 9 $\rightarrow$ 1 \\
0064680 & Epigallocatechin gallate & C$_{22}$H$_{18}$O$_{11}$ & 33 & 6 $\rightarrow$ 1 \\
0014764 & Verapamil & C$_{27}$H$_{38}$N$_2$O$_4$ & 33 & 1 \\
0028493 & Candesartan & C$_{24}$H$_{20}$N$_6$O$_3$ & 33 & 1 \\
0142445 & Diphenoxylate & C$_{30}$H$_{32}$N$_2$O$_2$ & 34 & 1 \\
0076066 & Letrozole & C$_{17}$H$_{11}$N$_5$ & 22 & 3 $\rightarrow$ 1 \\
0181733 & Unnamed MassSpecGym entry & C$_{16}$H$_{23}$N$_5$O & 22 & 1 \\
\bottomrule
\end{tabular}
\end{table}

\begin{table}[!htbp]
\centering
\caption{\textbf{Second-round candidate-guided W/V refinement examples.}
Local formula-preserving changes are proposed around high-ranked G candidates,
then evaluated using W/V evidence and chemical-topology filtering.}
\label{tab:second_round_refinement}
\small
\setlength{\tabcolsep}{3pt}
\begin{tabular}{@{}p{0.21\linewidth}p{0.29\linewidth}p{0.42\linewidth}@{}}
\toprule
Molecule & Initial candidate & Refinement and outcome \\
\midrule
Cyclic tetrapeptide (0002114) & Top-2; linear $n$-propyl side chain; Tanimoto 0.833 & Peripheral $n$-propyl $\rightarrow$ isopropyl reattachment; Exact GT recovered \\
Tetracycline (0058046) & Top-1; shifted keto--enol tautomer; Tanimoto 0.697 & Enumerate up to 64 tautomer states and score with W/V; Exact GT recovered \\
Epigallocatechin gallate (0064680) & Gen 6 Top-1 quinoid artifact; Tanimoto 0.431 & Reject with topology prior (score $-0.80$); promote previous rank 6; Exact GT promoted to rank 1 \\
\bottomrule
\end{tabular}
\end{table}

The Gen 7 candidate consensus stage uses Morgan-fingerprint centrality as a
post-generation ranking feature, together with a hand-specified chemical-topology
prior. Accordingly, the Gen 7 landmark results describe this particular
candidate-consensus ranking configuration and should not be described as a
fingerprint-free G-only result. The macro-core route also retrieves non-TEST
analogue structures, prunes them to strict partial graphs, and requires G to
complete the remaining atoms; it is therefore macro-core-assisted neural
completion rather than generation from an empty graph.

\subsection{Self-generated data and chemical-knowledge expansion}

\begin{table}[!htbp]
\centering
\caption{\textbf{Data and symbolic-knowledge resources used or targeted by the RSI programme.}
The status column distinguishes assets already generated and used in experiments
from external resources and code paths whose full integration remains in progress.}
\label{tab:rsi_data_assets}
\small
\setlength{\tabcolsep}{3pt}

\begin{tabular}{@{}p{0.25\linewidth}rp{0.32\linewidth}p{0.19\linewidth}@{}}
\toprule
Resource & Scale & Purpose & Status \\
\midrule
Non-VAL/TEST molecular structures & 597,383 molecules & Unlabelled structure pool for sleep rollouts and non-TEST macro-core/scaffold support & Source pool reported in RSI plan \\
GeMS-A de-leaked experimental MS/MS & 41,578,985 spectra & Real unlabeled spectra for wake-stream candidate inference & Resource identified; extraction script added \\
HMDB-W simulated MS/MS & 129,146 spectra & W-generated structure--spectrum pairs for G/MMA training & Generated and used in joint training \\
HMDB-W non-overlapping groups & 32,229 groups & Structure source for the HMDB-W simulation corpus and MMA expansion & Generated and used \\
MMA library v1 & 4,230,707 entries; 20,430 groups & TRAIN-derived fragment and motif vocabulary & Built \\
Unified MMA library v2 & 12,983,158 entries; 115,383 groups & TRAIN + HMDB-W fragment and motif vocabulary & Built; 220,642 ion-formula keys \\
\bottomrule
\end{tabular}
\end{table}

The new Dual Wake--Sleep script implements extraction of de-leaked GeMS queries
and construction of two feedback products: high-confidence real-spectrum
pseudo-labels and W-simulated hindsight pairs. The current script's default
query extraction is a bounded sampling configuration; PR95 does not yet report
a completed full-scale GeMS wake--sleep training cycle or a multi-generation
receipt demonstrating that all emitted pairs have been consumed by the next
generation. The pseudo-label gate uses explained intensity, W/V scores, a
positive candidate margin and a topology filter; hindsight pairs require
chemically valid, scaffold-diverse candidates with a sufficiently populated
W-emitted spectrum.

\begin{table}[!htbp]
\centering
\caption{\textbf{HMDB-W joint-training progression on held-out Calibration.}
The reported evaluation uses empty starts without W/V reranking or MMA starts;
the epoch trajectory therefore isolates the training progression more directly
than the full inference pipeline.}
\label{tab:hmdbw_epoch_progression}
\small
\setlength{\tabcolsep}{3pt}

\begin{tabular}{@{}rrrrrr@{}}
\toprule
Epoch & Train loss & Canonical@1 & Valid@1 & Exact@1 & Reach \\
\midrule
0 & 1.997 & 54.4\% & 64.7\% & 4.21\% & 7.37\% \\
2 & 1.690 & 58.4\% & 67.8\% & 7.37\% & 9.47\% \\
6 & 1.588 & 60.3\% & 69.6\% & \textbf{8.42\%} & \textbf{13.68\%} \\
9 & 1.575 & \textbf{60.6\%} & \textbf{69.8\%} & 7.37\% & 11.58\% \\
\bottomrule
\end{tabular}
\end{table}

\subsection{Failure diagnosis and fragment-level evidence}

\begin{table}[!htbp]
\centering
\caption{\textbf{Failure modes and action ambiguity in Peak-Ladder generation.}
The trajectory audit identifies common failure categories and quantifies the
number of competing valid actions under random atom ordering.}
\label{tab:ladder_failure_modes}
\small
\begin{tabular}{lr}
\toprule
Diagnostic & Result \\
\midrule
Wrong anchor attachment (\texttt{grow:anchor}) & 42\% of recorded errors \\
Wrong peak selection (\texttt{select:wrong\_peak}) & 21\% of recorded errors \\
Alternative growth action (\texttt{grow:add\_other}) & 10\% of recorded errors \\
States with more than one valid action & 46.8\% \\
Mean number of valid actions per state & 1.82 \\
Maximum number of valid actions in a state & 27 \\
Mean grow-stage ambiguity in failed trajectories & 2.17 \\
Mean grow-stage ambiguity in successful trajectories & 1.56 \\
\bottomrule
\end{tabular}
\end{table}

\begin{table}[!htbp]
\centering
\caption{\textbf{Strict fragment-tree combination trade-off.}
Combining more fragment anchors increases the fraction of reference atoms
covered, but decreases strict structural-hit accuracy under the corrected
substructure-matching definition.}
\label{tab:strict_fragment_tree_merge}
\small
\setlength{\tabcolsep}{3pt}
\begin{tabular}{@{}p{0.32\linewidth}rrrr@{}}
\toprule
Selection strategy & Covered atoms & GT coverage & Hit@0.8 & Coverage if hit \\
\midrule
Independent MMA Top-1 & 13.73 & 35.03\% & \textbf{46.09\%} & 19.90\% \\
Strict fragment-tree root & 15.88 & 38.99\% & 36.72\% & 19.81\% \\
Six-anchor merge & 18.64 & \textbf{40.86\%} & 22.66\% & 14.27\% \\
\bottomrule
\end{tabular}
\end{table}

\begin{table}[!htbp]
\centering
\caption{\textbf{MMA motif evidence complementary to the World fragment set.}
The union adds motif-supported fragments to the W set. The low complementary
win rates indicate that additional motif coverage does not automatically
provide discriminative evidence beyond W.}
\label{tab:w_mma_complementarity}
\small
\setlength{\tabcolsep}{3pt}
\begin{tabular}{@{}p{0.27\linewidth}rrrr@{}}
\toprule
Evidence set & Near-isomer win & Near margin & Random win & Random margin \\
\midrule
World W fragments & 81.20\% & 0.2674 & 87.36\% & 0.3897 \\
MMA motifs & 79.64\% & 0.2051 & 83.68\% & 0.2948 \\
MMA complement to W & 39.51\% & $-0.0450$ & 38.64\% & $-0.0543$ \\
W + MMA union & 80.17\% & 0.2224 & 87.05\% & 0.3353 \\
\bottomrule
\end{tabular}
\end{table}

In the same audit, W left 21.61\% of the reference fragment intensity unexplained;
MMA motifs covered approximately 15.00\% of that W-unexplained intensity. The
union did not improve the near-isomer win rate over W alone. This result motivates
joint structural compatibility modelling rather than treating every additional
motif match as independent positive evidence.

\begin{table}[!htbp]
\centering
\caption{\textbf{Early MMA-conditioned generation diagnostic.}
The SEARCH\_DEV comparison shows improvements in structural similarity and
candidate-pool quality from MMA conditioning, while exact connectivity recovery
remains low. These results are from an earlier generator and are not directly
comparable with the later held-out Calibration or official-test evaluations.}
\label{tab:early_mma_search_dev}
\small
\setlength{\tabcolsep}{3pt}
\begin{tabular}{@{}p{0.32\linewidth}rrrrr@{}}
\toprule
Generator configuration & Exact@1 & Exact@10 & Reach@100 & Tan@1 & Top-10 max Tan \\
\midrule
Baseline generator & 0.78\% & 1.56\% & 1.56\% & 0.1424 & 0.1951 \\
Baseline + MMA & 0.00\% & 1.56\% & 1.56\% & 0.1672 & \textbf{0.2406} \\
Baseline + mixed motifs & 0.00\% & 1.56\% & 1.56\% & 0.1611 & 0.2252 \\
\bottomrule
\end{tabular}
\end{table}

\paragraph{RSI implementation status.}
PR95 includes a four-phase RSI controller for self-play and failure mining,
symbolic MMA/Ring vocabulary expansion, Hybrid G training, and multi-gate
promotion or rollback. It also includes a feedback miner that extracts
same-formula isomer decoys and newly covered ring systems from G rollouts.
These components provide the execution framework for recursive updates.
The Gen~1--7 ledger above records successive experimental configurations;
a complete multi-generation autonomous ledger, with every generation's
hypothesis, feedback products, model promotion and subsequent performance,
is not yet reported. Likewise, the Dual Wake--Sleep code is implemented as a
data-generation path, but its full end-to-end training cycle remains pending.

\paragraph{Source and provenance.}
The tables in this section are transcribed from the PR95 progress report and
the linked machine-readable PR93 artifacts, including
\texttt{pr93\_test\_generalization\_and\_mma\_ablation\_20261109.md},
\texttt{pr93\_rsi\_evolution\_plan\_and\_gen7\_report\_20261109.md},
\texttt{v\_mma\_complement\_audit.json},
\texttt{anchor\_strict\_merge1.json},
\texttt{anchor\_strict\_merge6.json}, and
\texttt{compare\_pmma3.json}. The source implementation and experiment receipts
are linked from \url{https://github.com/wangyu-sd/orbit-ms/pull/95}.

% Flush supplementary table floats before the bibliography.
\clearpage
